\section{Trimming across one cut}

Landau, Vazirani and Vidick~\cite{Landau2015Polynomial} (arXiv:1307.5143v1,
lines 330--333) trim a chain vector across a cut by keeping the Schmidt terms
with the largest coefficients. Trimming across one cut does not increase the
Schmidt rank across any other cut (lines 544--553).

\begin{definition}[Trimming a chain vector]
    \label{def:chain_trim}
    \lean{TNLean.MPS.AreaLaw.chainTrim}
    \leanok
    \uses{def:chain_open_mps}
    Let $\psi$ be a vector on $L$ sites of local dimension $q$, let $i$ and
    $m$ be natural numbers, and write $C=M_i(\psi)$ for the coefficient matrix
    at the cut $i$. A \emph{trimming projection} of $\psi$ at $m$ across the
    cut $i$ is an orthogonal projection $P$ on $\mathcal H_{[0,i)}$ of rank at
    most $m$ with $\operatorname{tr}(PCC^{*})$ equal to the sum of the $m$
    largest eigenvalues of $CC^{*}$. The \emph{trimmed vector} is
    $\operatorname{trim}^i_P\psi=(P\otimes1)\psi$, whose coefficient matrix
    at the cut $i$ is $PC$. When a Schmidt coefficient is repeated, every
    choice of kept eigenvectors gives a trimming projection.
\end{definition}

\begin{theorem}[Trimming does not increase cut ranks]
    \label{thm:chain_trim_cut_ranks}
    \lean{
        TNLean.MPS.AreaLaw.cutMatrix_chainTrim,
        TNLean.MPS.AreaLaw.cutVector_chainTrim,
        TNLean.MPS.AreaLaw.rank_cutMatrix_chainTrim_self_le,
        TNLean.MPS.AreaLaw.rank_cutMatrix_chainTrim_le
    }
    \leanok
    \uses{def:chain_trim}
    Let $P$ be a trimming projection of $\psi$ at $m$ across the cut $i$.
    The coefficient matrix of $\operatorname{trim}^i_P\psi$ at the cut $i$ is
    $PM_i(\psi)$, so $\operatorname{rank}M_i(\operatorname{trim}^i_P\psi)\le m$, and
    $\operatorname{rank}M_j(\operatorname{trim}^i_P\psi)\le
        \operatorname{rank}M_j(\psi)$ at every cut $j$. No assumption on $L$
    or $q$ is needed.
\end{theorem}
\begin{proof}\leanok
    \uses{thm:hastings_supported_cut_rank}
    The first bound holds because $PC$ has rank at most that of $P$. For
    $j\ge i$ the operator $P\otimes1$ acts on sites to the left of the cut
    $j$. For $j\le i$, the projection $P$ commutes with $CC^{*}$, since it
    maximizes $\operatorname{tr}(RCC^{*})$ among orthogonal projections $R$
    of the same trace. Hence $PC=CQ^{\mathsf T}$, where $Q^{\mathsf T}$ is the
    support projection of $C^{*}PC$, so $(P\otimes1)\psi=(1\otimes Q)\psi$
    with $Q$ acting on sites to the right of the cut $j$. In both cases the
    operator acts on one side of the cut $j$ and cannot increase its rank.
    If there are no configurations, both coefficient matrices are empty.
\end{proof}

\begin{theorem}[Trimming error on the chain]
    \label{thm:chain_trim_error}
    \lean{
        TNLean.MPS.AreaLaw.norm_chainTrim_le,
        TNLean.MPS.AreaLaw.norm_inner_sub_chainTrim_le
    }
    \leanok
    \uses{def:chain_trim}
    Let $P$ be a trimming projection of $\psi$ at $m$ across the cut $i$,
    and let $w$ be a vector with $\operatorname{rank}M_i(w)\le D$. Then
    $\|\operatorname{trim}^i_P\psi\|\le\|\psi\|$ and
    \begin{align}
        |\langle\psi-\operatorname{trim}^i_P\psi,w\rangle|
        \le\|\psi\|\,\|w\|\sqrt{\frac{D}{m+1}}. \notag
    \end{align}
    The source states error $\epsilon$ for $m=D/\epsilon$, which is false;
    the bound above gives error $\epsilon$ for $m+1\ge D/\epsilon^2$
    \cite{gap:lvv13_trimming_radius}.
\end{theorem}
\begin{proof}\leanok
    The cut coordinates at $i$ preserve norms and inner products, and in
    them the trimmed vector is $(P\otimes1)$ applied to the cut vector of
    $\psi$. The squared norm of $(P\otimes1)\psi$ is the sum of the $m$
    largest eigenvalues of $CC^{*}$, at most their total $\|\psi\|^2$. For the
    second bound, let $T$ be the support projection of the part of $M_i(w)$
    orthogonal to the range of $P$. Then $P+T$ is an orthogonal projection of
    rank at most $m+D$, so by Ky Fan's maximum principle the discarded part
    paired with $w$ captures at most $D$ times the $(m+1)$-st eigenvalue of
    $CC^{*}$, which is at most $\|\psi\|^2/(m+1)$.
\end{proof}

\subsection{Constant rank across one cut}

Landau, Vazirani and Vidick~\cite{Landau2015Polynomial} (lines 342--364)
use that the ground state of a gapped chain can be trimmed across one cut to
a bond dimension independent of the length of the chain, and that trimming a
vector of low energy keeps its energy low. The bond dimension they state,
$\exp(O(\varepsilon^{-1}\log^3d\log(1/\delta)))$, comes from Arad, Kitaev,
Landau and Vazirani~\cite{Arad2013AreaLaw}; the one below comes from the
area law of this chapter and has different constants
\cite{gap:lvv13_constant_bond}. The energy bound after trimming is the
corrected one of \cite{gap:lvv13_trimming_energy}.

\begin{definition}[A constant bond dimension]
    \label{def:chain_constant_bond}
    \lean{TNLean.MPS.AreaLaw.constantCutBond}
    \leanok
    \uses{def:chain_model}
    For a chain of local dimension $q$ with range $R$, term norms at most $J$
    and gap $\Delta>0$, let $C_{\mathrm{int}}$ be the interval entropy bound,
    $s=\sqrt{(1+J/\Delta)B_R}$, and $0<\delta\le1$. Put
    \begin{align}
        B_\delta=\Bigl\lfloor\exp\Bigl(C_{\mathrm{int}}
        +32s\log\frac{2e^{e/2}}{\delta}\Bigr)\Bigr\rfloor. \notag
    \end{align}
    It depends neither on the length of the chain nor on a cut.
\end{definition}

\begin{theorem}[Trimming the ground state to a constant bond dimension]
    \label{thm:chain_constant_bond_overlap}
    \lean{
        TNLean.MPS.AreaLaw.ChainHamiltonian.one_sub_le_re_inner_chainTrim,
        TNLean.MPS.AreaLaw.inner_normalized_chainTrim,
        TNLean.MPS.AreaLaw.cutVector_normalized_chainTrim,
        TNLean.MPS.AreaLaw.regionState_leftBlock_eq_leftMarginal
    }
    \leanok
    \uses{def:chain_trim, def:chain_constant_bond}
    Let $\Omega$ be a gapped ground vector, with gap $\Delta>0$, of a
    finite-range Hamiltonian with term norms at most $J\ge0$ on a chain of
    local dimension $q\ge1$, and let $0<\delta\le1$. For every cut $k$ and
    every trimming projection $P$ of $\Omega$ at $B_\delta$ across the cut $k$,
    the vector $v=\operatorname{trim}^k_P\Omega/
        \|\operatorname{trim}^k_P\Omega\|$ satisfies
    $\operatorname{Re}\langle v,\Omega\rangle\ge1-\delta$, and the overlap is
    real and nonnegative. This is the statement of the source with the bond
    dimension of Definition~\ref{def:chain_constant_bond} in place of the
    stated one \cite{gap:lvv13_constant_bond}.
\end{theorem}
\begin{proof}\leanok
    \uses{thm:chain_prefix_spectral_bounds}
    The reduced state $\rho$ of $\Omega$ on the first $k$ sites is $CC^{*}$
    for the coefficient matrix $C$ at the cut. Take
    $w=32s\log(2e^{e/2}/\delta)$ in
    Theorem~\ref{thm:chain_prefix_spectral_bounds}: the spectral projection
    there has rank at most $B_\delta$ and discards mass at most $\delta$. By
    Ky Fan's maximum principle the sum of the $B_\delta$ largest eigenvalues
    of $\rho$ is at least $1-\delta$. This sum is
    $\|\operatorname{trim}^k_P\Omega\|^2$, and
    $\langle v,\Omega\rangle=\|\operatorname{trim}^k_P\Omega\|\ge
        \sqrt{1-\delta}\ge1-\delta$.
\end{proof}

\begin{theorem}[The energy after trimming]
    \label{thm:chain_trim_energy}
    \lean{
        TNLean.MPS.AreaLaw.re_inner_normalized_chainTrim_le,
        TNLean.MPS.AreaLaw.ChainHamiltonian.re_inner_normalized_chainTrim_le,
        TNLean.MPS.AreaLaw.rank_cutMatrix_normalized_chainTrim_le,
        TNLean.MPS.AreaLaw.norm_normalized_chainTrim_eq_one,
        Matrix.IsGappedGroundState.posSemidef_sub_smul_one
    }
    \leanok
    \uses{def:chain_trim, thm:chain_constant_bond_overlap}
    Let $H=H_L+H_R+K$ on a chain, with $H_L$ supported on the first $k$
    sites, $H_R$ on the remaining sites and $0\le K\le1$, and let $\Omega$
    be a gapped ground vector of $H$ with energy $E_0$ and gap $\Delta>0$.
    Let $B$ be a bond dimension and $P_\Omega$ a trimming projection of
    $\Omega$ at $B$ across the cut $k$ whose normalized trimmed vector $u$
    satisfies $\operatorname{Re}\langle u,\Omega\rangle\ge1-\delta_B$. Let
    $w$ be a unit vector with $\operatorname{Re}\langle w,Hw\rangle\le
        E_0+\delta$, and $P$ a trimming projection of $w$ at $B$ across the cut
    $k$. If $\eta\ge2(\delta/\Delta+\delta_B)$ and $2\eta<1$, the normalized
    trimmed vector $v$ of $w$ satisfies
    \begin{align}
        \operatorname{Re}\langle v,Hv\rangle
        \le E_0+\frac{\delta+2\sqrt{2\eta}}{1-2\eta}. \notag
    \end{align}
    For a finite-range chain Hamiltonian that splits in this way at the cut
    $k$, this holds at $B=B_{\delta_B}$ for every $0<\delta_B\le1$. In the
    source $H$ is a sum of nearest-neighbour terms $0\le H_i\le1$ and $K$ is
    the term across the cut; its printed bound $E_0+6\sqrt\delta$ does not
    follow from its argument \cite{gap:lvv13_trimming_energy}.
\end{theorem}
\begin{proof}\leanok
    The gap gives $H\ge E_0$, so $\delta\ge0$, and
    $|\langle\Omega,w\rangle|\ge1-\delta/\Delta$. Since $u$, $w$ and $\Omega$
    are unit vectors, the overlap estimate for three unit vectors gives
    $|\langle u,w\rangle|\ge1-2(\delta_B+\delta/\Delta)\ge1-\eta$. In the
    coordinates of the cut $k$, which preserve norms and inner products,
    $H_L$ becomes $A\otimes1$, $H_R$ becomes $1\otimes B'$, and $u$ has
    Schmidt rank at most $B$. The bipartite energy estimate for trimmed
    vectors then gives the bound: the trimming projection commutes with the
    left marginal of $w$, so $A\otimes1+1\otimes B'$ has no cross terms
    between the trimmed part and the discarded part, the cross term of $K$ is
    at most twice the norm of the discarded part, and by the Eckart--Young
    theorem the squared norm of the discarded part is at most $2\eta$.
\end{proof}

\begin{theorem}[A state of constant rank across one cut]
    \label{thm:chain_constant_rank_state}
    \lean{TNLean.MPS.AreaLaw.ChainHamiltonian.exists_rank_cutMatrix_le_constantCutBond}
    \leanok
    \uses{def:chain_constant_bond}
    Let $\Omega$ be a gapped ground vector, with energy $E_0$ and gap
    $\Delta>0$, of a finite-range chain Hamiltonian with term norms at most
    $J\ge0$ and local dimension $q\ge1$, which splits at the cut $k$ as in
    Theorem~\ref{thm:chain_trim_energy}, and let $0<\delta<1/4$. There is a
    unit vector $v$ with $\operatorname{rank}M_k(v)\le B_\delta$,
    $\operatorname{Re}\langle v,\Omega\rangle\ge1-\delta$, and
    \begin{align}
        \operatorname{Re}\langle v,Hv\rangle\le
        E_0+\frac{4\sqrt\delta}{1-4\delta}. \notag
    \end{align}
    The source states the energy bound $E_0+6\sqrt\delta$
    \cite{gap:lvv13_trimming_energy}.
\end{theorem}
\begin{proof}\leanok
    \uses{thm:chain_constant_bond_overlap, thm:chain_trim_energy}
    Take $v$ the normalized trimmed vector of $\Omega$ at $B_\delta$. Its
    overlap is Theorem~\ref{thm:chain_constant_bond_overlap}, which also
    shows that the trimmed vector is nonzero, and its rank across the cut is
    at most $B_\delta$. Theorem~\ref{thm:chain_trim_energy} with $w=\Omega$,
    energy excess $0$, $\delta_B=\delta$ and $\eta=2\delta$ gives
    $E_0+2\sqrt{4\delta}/(1-4\delta)$.
\end{proof}
